\Needspace{18\baselineskip}
\section{Finite interaction patterns and simple-context search}
\label{sec:contextual-search}
This section discharges the search adequacy obligation formerly listed in
Appendix~\ref{app:remaining-work}. The principal theorem identifies the checked
erasure-based search relation with actual plugging of finite interaction
fragments into one focus-free simple hole context, in both directions. It also
identifies the independent incidence presentation with that same relation.
The result concerns search: it does not add substitution or compute transitions.

We use the five coordinate clauses already specified for Fig.~9 to index pairs
of actual two-edge graphs. Each graph contains one variable or operation edge
and one focus edge. A single shallow slot supplies the surrounding graph;
opening it creates an explicitly labelled hole. General slot-label transport
and a closing theorem prove that arbitrary simple contexts are covered, rather
than restricting context application to a chosen family of examples.

\subsection{Subdividing at an output port}
\begin{definition}[Incoming-side focus subdivision]\label{term:output-focus-insertion}
For a graph $G:I\to O$ and a term port $v$ that is not an input, retain the old
carriers and add one internal term vertex $z$ and one focus edge $f$. Original
source lists are mapped through the old-vertex inclusion. Original target lists
are mapped through the split function that replaces $v$ by $z$ and includes
all other ports. Set the focus source to $z$ and its target to the included
copy of $v$. Thus a focus can be placed at an output boundary, which the earlier
outgoing-side subdivision cannot do.
\end{definition}
\lean{Graph/OutputFocusInsertion}{9}{11}{Linear.OutputFocusInsertion.target}
\lean{Graph/OutputFocusInsertion}{13}{15}{Linear.OutputFocusInsertion.source}
\begin{lemma}[Incidence uniqueness and coverage]
The incoming subdivision satisfies source and target uniqueness and the two
coverage laws of a finite graph.
\end{lemma}
\begin{proof}
For two original target lists, injectivity of splitting reflects shared
membership to the original graph's target uniqueness. An original target
cannot equal the focus target: a split port never equals the unsplit inclusion
of the selected port. Two focus-edge names coincide because their carrier is a
singleton. For sources, old-vertex inclusion is injective and its image is
disjoint from the fresh focus source. Original source uniqueness and the
singleton focus carrier then give source uniqueness.

For target coverage, a retained port unequal to $v$ has its original target
edge, a retained copy of $v$ is targeted by the focus, and the fresh vertex is
targeted by the original incoming edge of $v$. That edge exists because $v$
is not an input. Conversely each target is not an input, using the original
coverage law, the hypothesis on $v$, or the internal status of $z$. Source
coverage retains every old source incidence and supplies the fresh vertex
with the focus edge; no input-side exception is needed for this direction.
\end{proof}
\lean{Graph/OutputFocusInsertion}{17}{36}{Linear.OutputFocusInsertion.target_unique}
\lean{Graph/OutputFocusInsertion}{38}{57}{Linear.OutputFocusInsertion.source_unique}
\lean{Graph/OutputFocusInsertion}{59}{86}{Linear.OutputFocusInsertion.target_coverage}
\lean{Graph/OutputFocusInsertion}{88}{106}{Linear.OutputFocusInsertion.source_coverage}
The constructor packages these incidence functions and laws as a typed finite
hypernet. Finite enumerations are disjoint sums with a singleton. Nodup follows
from injectivity of the host or split map. The existing host-label and split-label
lemmas prove typing of original incidences; the new focus has a singleton term
source and target. Given global focus freedom of $G$, the focus certificate
selects the fresh edge, retains freedom of every old label, and proves distinct
endpoints by freshness. Its source is internal, so it cannot be exposed from
an input directly to an output.
\lean{Graph/OutputFocusInsertion}{109}{158}{Linear.OutputFocusInsertion.insert}
\lean{Graph/OutputFocusInsertion}{161}{177}{Linear.OutputFocusInsertion.focussed}

\subsection{Finite fragments and their plugging equations}
\begin{definition}[Boundary focus fragment]\label{term:search-fragment}
For a selected edge $e$ of $C$, its corolla has the same label and ordered
interface, no internal vertices, and one edge. The input fragment at a term
input index inserts a focus before that corolla input. The output fragment at
a term output index uses incoming subdivision after that corolla output.
Both fragments have one internal vertex and two edges. The corresponding
whole-graph subdivisions insert at the selected edge's source or target port.
\end{definition}
\lean{Graph/PatternPlugging}{10}{10}{Linear.PatternPlugging.corolla}
\lean{Graph/PatternPlugging}{42}{43}{Linear.PatternPlugging.inputPattern}
\lean{Graph/PatternPlugging}{46}{47}{Linear.PatternPlugging.inputWhole}
\lean{Graph/PatternPlugging}{118}{119}{Linear.PatternPlugging.outputPattern}
\lean{Graph/PatternPlugging}{122}{123}{Linear.PatternPlugging.outputWhole}
The vertex bijection from a plugged fragment to the corresponding whole
subdivision removes the empty corolla-vertex summand and retains the host and
fresh vertices. The edge bijection retains every unselected host edge, restores
the selected edge name from the corolla edge, and sends the fragment focus to
the fresh whole-graph focus. Its inverse tests whether an old edge is the
selected edge; the singleton carriers prove the inverse equations.
\lean{Graph/PatternPlugging}{13}{17}{Linear.PatternPlugging.vertices}
\lean{Graph/PatternPlugging}{20}{39}{Linear.PatternPlugging.edges}

\begin{lemma}[Input-fragment incidence equations]
Under those bijections, host ports retain their host inclusions; included
corolla ports become inclusions of their attachment ports; and a split corolla
input becomes the split attached source port.
\end{lemma}
\begin{proof}
The host equation is a case analysis on internal, input, and output ports.
For an included corolla port, the internal case is empty; the input and output
cases reduce to that host equation at the corresponding attachment. For a
split input, test equality of its index with the selected index. If equal,
both sides are the fresh vertex. Otherwise injectivity of ordered source-port
lookup proves that the attachment differs from the selected attachment, so
both sides are the retained host port.
\end{proof}
\lean{Graph/PatternPlugging}{49}{51}{Linear.PatternPlugging.input_host}
\lean{Graph/PatternPlugging}{53}{61}{Linear.PatternPlugging.input_insert_host}
\lean{Graph/PatternPlugging}{63}{73}{Linear.PatternPlugging.input_insert_split}
\begin{theorem}[Input-fragment factorization]
Plugging the input fragment at $e$ is isomorphic to subdividing the whole graph
at that edge's selected source port.
\end{theorem}
\begin{proof}
Use the carrier bijections above. Labels are unchanged. A source port of an
unselected host edge cannot equal the selected source port: otherwise source
uniqueness would identify its edge with $e$. Its source list therefore agrees
with the whole-graph split list. For the selected edge, enumerate the corolla's
input indices and apply the split-input equation pointwise; enumerating the
host's source ports reconstructs its source list. The focus source is the
included attachment. All old targets use the host equation, and enumerating
corolla outputs reconstructs the selected target list. Both focus targets are
fresh. These cases establish every isomorphism field.
\end{proof}
\lean{Graph/PatternPlugging}{76}{115}{Linear.PatternPlugging.input_iso}

The output-fragment equations are proved separately with targets in place of
sources. Host and included-port equations again follow by port cases. For
split outputs, equality of output indices selects the fresh vertex; inequality
is reflected by injectivity of ordered target-port lookup.
\lean{Graph/PatternPlugging}{125}{127}{Linear.PatternPlugging.output_host}
\lean{Graph/PatternPlugging}{129}{137}{Linear.PatternPlugging.output_insert_host}
\lean{Graph/PatternPlugging}{139}{149}{Linear.PatternPlugging.output_insert_split}
\begin{theorem}[Output-fragment factorization]
Plugging the output fragment at $e$ is isomorphic to incoming subdivision of
the whole graph at its selected target port.
\end{theorem}
\begin{proof}
Use the same carrier bijections. Target uniqueness ensures that an unselected
host edge never targets the selected target port, so its target list is
unchanged by splitting. Enumerate output indices and apply the split-output
equation for the selected edge. Its focus target is the included attachment.
Sources use host inclusion, with the selected source list reconstructed by
its ordered input enumeration; the focus source is fresh. Labels and finite
carriers are preserved exactly as for input fragments.
\end{proof}
\lean{Graph/PatternPlugging}{152}{185}{Linear.PatternPlugging.output_iso}
Neither factorization assumes that an input attachment differs from an output
attachment. Each uses injectivity only within the relevant ordered side.
This distinction is necessary for self-looping slots.

\subsection{Reconciling the two subdivision orientations}
\begin{definition}[Subdivision vertex exchange]\label{term:subdivision-exchange}
At an internal port $v$, exchange its old copy with the fresh focus vertex and
fix every other old vertex. This self-inverse function is a bijection on the
extended internal carrier; singleton uniqueness handles the fresh carrier's
only index.
\end{definition}
\lean{Graph/FocusSubdivision}{10}{20}{Linear.FocusSubdivision.swap}
Its port extension sends an included old port to its split version and the
fresh port to the included copy of $v$. The first claim follows by cases on
the port and equality with $v$; the second follows directly from the definition.
Applying the exchange again shows that a split port returns to its inclusion.
\lean{Graph/FocusSubdivision}{22}{28}{Linear.FocusSubdivision.swap_host}
\lean{Graph/FocusSubdivision}{30}{30}{Linear.FocusSubdivision.swap_fresh}
\lean{Graph/FocusSubdivision}{32}{35}{Linear.FocusSubdivision.swap_split}
\begin{theorem}[Agreement of subdivision orientations]
Incoming and outgoing focus subdivision at an internal term port are
isomorphic, with the focus edge and its mode retained.
\end{theorem}
\begin{proof}
Use the exchange on vertices and the identity on edges. The two exchanged
vertices have term type; all other labels are fixed. On original sources,
the exchange turns inclusion into splitting. On original targets it turns
splitting into inclusion. Its fresh-port and selected-port equations identify
the focus source and target lists. Edge labels are unchanged, including the
focus mode.
\end{proof}
\lean{Graph/FocusSubdivision}{38}{68}{Linear.FocusSubdivision.iso}

\subsection{Factoring the search rules through one corolla}
\begin{definition}[Corolla attachment]\label{term:corolla-attachment}
Attach each corolla input index to the corresponding source port of $e$, and
each output index to the corresponding target port. There is no internal case.
Typing of ordered port lookup proves preservation of vertex labels. Mapping
attachment over either corolla incidence list reconstructs the corresponding
host list by ordered port enumeration.
\end{definition}
\lean{Graph/CorollaSearch}{10}{13}{Linear.PatternPlugging.attach}
\lean{Graph/CorollaSearch}{15}{20}{Linear.PatternPlugging.attach_label}
\lean{Graph/CorollaSearch}{22}{25}{Linear.PatternPlugging.attach_source}
\lean{Graph/CorollaSearch}{27}{30}{Linear.PatternPlugging.attach_target}
An operation view on the corolla becomes a view on $e$ by attaching its source,
argument list, and thunk list. The reconstructed incidence equations and
preservation of list lengths prove the view's fields. Attachment commutes
with entry by inspecting whether the argument list is empty, and with return
by commuting lookup with mapping and inspecting the next argument.
\lean{Graph/CorollaSearch}{33}{42}{Linear.PatternPlugging.operation}
\lean{Graph/CorollaSearch}{44}{46}{Linear.PatternPlugging.operation_enter}
\lean{Graph/CorollaSearch}{48}{51}{Linear.PatternPlugging.operation_resume}
\begin{lemma}[Rule attachment and operation recovery]
Every coordinate rule on the corolla attaches to a coordinate rule of $C$.
Every operation view of $C$ is recovered from a view on its edge's corolla,
with equal source, ordered arguments, and operation label after attachment.
\end{lemma}
\begin{proof}
For a variable rule, reconstruct the singleton source list by attachment and
apply the host variable rule. For entry, use the attached operation view and
entry commutation. For return, map the argument lookup and use return
commutation. For operation recovery, the corolla has the original operation
label and hence admits a view. Its attached view has the original edge name;
view uniqueness identifies the required data.
\end{proof}
\lean{Graph/CorollaSearch}{54}{65}{Linear.PatternPlugging.sound}
\lean{Graph/CorollaSearch}{68}{75}{Linear.PatternPlugging.operation_preimage}
\begin{theorem}[Finite-rule factorization]
Every coordinate search step on $C$ is the attachment of a search step on the
corolla of one interacting edge, with both endpoint coordinates attached by
the same map.
\end{theorem}
\begin{proof}
Analyze the search rule. For a variable, the signature forces its corolla
source list to have length one. Extract that source port; reconstruction of
the host source list identifies its attachment with the original source.
The corolla variable rule has the required two endpoints.
For entry, recover the corolla operation view. Equality of source data and
uniqueness of entry selection give the endpoint equations. For return,
recover the view and invert lookup in its mapped argument list. A failed
lookup contradicts the given host lookup. A successful lookup supplies the
same argument index and a port attaching to the host argument. The corolla
return rule and uniqueness of return selection give the target equation.
\end{proof}
\lean{Graph/CorollaSearch}{79}{112}{Linear.PatternPlugging.complete}
Inspecting any corolla rule shows that its label is either the variable label
or an operation label: the variable constructor supplies the first, and the
operation views in entry and return supply the second. Such a primitive label
is focus-free and has no hole in its occurrence fibre. At either depth form,
its shallow occurrence is not a hole and its child-occurrence type is empty.
These facts are derived from the rule, not assumed of context application.
\lean{Graph/CorollaSearch}{115}{122}{Linear.PatternPlugging.interaction_label}
\lean{Graph/CorollaSearch}{126}{137}{Linear.PatternPlugging.interaction_closed}

\subsection{Actual state endpoints}
A literal graph isomorphism between globally uniquely focussed graphs must
preserve their distinguished focuses. Otherwise the image of the source focus
would be a non-focus edge of the target and therefore have a focus-free label,
contradicting label preservation. Once the edge is identified, injectivity of
focus labels identifies its mode. This supplies a focussed isomorphism from
an ordinary graph isomorphism.
\lean{Graph/PatternRealization}{10}{27}{Linear.Focussed.isoOfGraph}
\begin{definition}[Typed fragment coordinate and matching]\label{term:pattern-matching}
A fragment coordinate is a mode and term port of the selected corolla.
Unlike a state coordinate it may lie at an output. Its fragment is the input
or output graph defined above, according to the port constructor. Attaching
it gives a typed state coordinate, by label preservation. A fragment matches
an actual state when its plugging is literally graph-isomorphic to that state.
\end{definition}
\lean{Graph/PatternRealization}{30}{30}{Linear.PatternPlugging.Position}
\lean{Graph/PatternRealization}{33}{39}{Linear.PatternPlugging.fragment}
\lean{Graph/PatternRealization}{42}{43}{Linear.PatternPlugging.attached}
\lean{Graph/PatternRealization}{78}{79}{Linear.PatternPlugging.Matches}
Every output attachment in a host of interface $\star\to[]$ is internal:
its target-list membership excludes inputs, and there are no output indices.
\lean{Graph/PatternRealization}{46}{51}{Linear.PatternPlugging.output_internal}
\begin{theorem}[Fragment realization]
In a focus-free, hole-free state-interface host, plugging any typed fragment
is isomorphic to realization of its attached coordinate. Consequently matching
an actual state is equivalent to a focussed isomorphism from that realization.
\end{theorem}
\begin{proof}
For input coordinates apply input-fragment factorization. For output
coordinates apply output-fragment factorization, then exchange the old
internal attachment and fresh vertex to obtain outgoing subdivision.
The previous internality lemma justifies that exchange. The corolla has no
internal coordinate case. Compose this graph isomorphism with a matching
isomorphism and use uniqueness of the focus to obtain the focussed endpoint
condition. Conversely discard the focus fields and compose graph
isomorphisms in the opposite direction.
\end{proof}
\lean{Graph/PatternRealization}{54}{75}{Linear.PatternPlugging.realization}
\lean{Graph/PatternRealization}{83}{90}{Linear.PatternPlugging.matches_iff}

\begin{definition}[Slot-context search]\label{term:slot-context-search}
Choose a focus-free, hole-free host, one selected edge, and two typed corolla
coordinates related by a finite search rule. Require their two plugged
fragments to match the source and target actual states. The same selected
edge is removed by both pluggings.
\end{definition}
\lean{Graph/ContextualSearch}{11}{14}{Linear.State.ContextualSearch}
\begin{theorem}[Slot-context correspondence]
Slot-context search and erasure-based state search coincide in both directions.
They therefore also coincide with incidence search. Recursive state search
adds precisely recursive renaming of a slot-context successor.
\end{theorem}
\begin{proof}
For soundness, attach the finite corolla rule to the host. Realize the resulting
coordinate step and transport its source and target through the matching
focussed isomorphisms. For completeness, erase the source state's focus and
factor its coordinate search through a corolla. The source coordinate is
typed by attachment-label preservation; search type preservation types the
target coordinate. The source endpoint matches by general focus recovery,
and the target matches by the endpoint isomorphism in the state search witness.
The incidence equivalence follows from the previous section's theorem.
Substitution into the characterization of recursive state search proves the
recursive endpoint statement.
\end{proof}
\lean{Graph/ContextualSearch}{17}{23}{Linear.State.contextualSearch_sound}
\lean{Graph/ContextualSearch}{27}{46}{Linear.State.contextualSearch_complete}
\lean{Graph/ContextualSearch}{49}{51}{Linear.State.search_iff_contextual}
\lean{Graph/ContextualSearch}{54}{56}{Linear.State.incidenceSearch_iff_contextual}
\lean{Graph/ContextualSearch}{59}{62}{Linear.State.deepSearch_iff_contextual}

\subsection{Opening and closing an explicit simple context}
\begin{definition}[Signature-preserving slot relabelling]\label{term:slot-relabelling}
Replace the selected edge label by a label with equal ordered input and output
signatures, keeping all other data. The selected-label and other-label equations
follow by the equality test defining relabelling. Identity carrier maps give
an incidence isomorphism with the unrestricted label relation. The supplied
signature equalities identify the old and new selected interfaces.
\end{definition}
\lean{Graph/SearchContexts}{16}{20}{Linear.SearchContexts.relabel}
\lean{Graph/SearchContexts}{22}{24}{Linear.SearchContexts.selected_label}
\lean{Graph/SearchContexts}{26}{29}{Linear.SearchContexts.other_label}
\lean{Graph/SearchContexts}{32}{38}{Linear.SearchContexts.geometry}
\lean{Graph/SearchContexts}{40}{42}{Linear.SearchContexts.input_eq}
\lean{Graph/SearchContexts}{44}{46}{Linear.SearchContexts.output_eq}
\begin{theorem}[The discarded slot label is immaterial]
Plugging any replacement before or after slot relabelling gives isomorphic
graphs, after transport along the equal interfaces. This works in both directions.
\end{theorem}
\begin{proof}
Apply the existing host-plugging transport to the identity incidence
isomorphism. For an unselected host edge the labels actually agree, and an
inserted edge retains its label, so the transported incidence isomorphism is
a literal graph isomorphism. For the reverse direction transport a replacement
back to the old interface, apply the forward result, and cancel consecutive
opposite interface transports. Cancellation follows by eliminating the two
interface equalities and simplifying identity transport. Invert the result.
\end{proof}
\lean{Graph/SearchContexts}{49}{65}{Linear.SearchContexts.plug_iso}
\lean{Graph/SearchContexts}{69}{80}{Linear.SearchContexts.plug_iso_reverse}
Relabelling by a focus-free label preserves global focus freedom: inspect the
selected edge separately and retain the other labels. If every existing hole
occurrence belongs to the selected edge, and the new label's entire fibre
has no hole, relabelling removes all holes. A hypothetical hole in the new
selected fibre contradicts its emptiness; one at another edge transports to
an old hole occurrence whose owner would have to be the selected edge.
\lean{Graph/SearchContexts}{83}{89}{Linear.SearchContexts.relabel_free}
\lean{Graph/SearchContexts}{93}{110}{Linear.SearchContexts.close_holeFree}

\begin{definition}[Opened slot]\label{term:opened-search-slot}
Give a hole any natural-number name and the selected edge's complete input
and output signatures. Opening the slot relabels that edge by this hole.
The selected label is exactly the specified hole label. If the host is
focus-free, so is the opened graph: a hole label contains no focus.
\end{definition}
\lean{Graph/SearchContexts}{113}{114}{Linear.SearchContexts.hole}
\lean{Graph/SearchContexts}{117}{119}{Linear.SearchContexts.openSlot}
\lean{Graph/SearchContexts}{121}{122}{Linear.SearchContexts.open_label}
\lean{Graph/SearchContexts}{125}{133}{Linear.SearchContexts.focusFree}
\begin{lemma}[Exactly one shallow hole]
Opening a slot in a hole-free host gives a context with exactly one shallow
hole occurrence, of the specified label, and no nested holes.
\end{lemma}
\begin{proof}
Any occurrence at another edge transports along its unchanged label to a hole
in the original host, contradicting hole freedom. The selected edge's fibre
is a shallow singleton plus an empty child type, and is therefore a
subsingleton. Two hole occurrences thus have the same owner and equal fibre
positions, proving global hole uniqueness. The shallow occurrence explicitly
witnesses presence. Finally any hole occurrence has that owner and that
shallow position, so its label is exactly the chosen hole.
\end{proof}
\lean{Graph/SearchContexts}{137}{148}{Linear.SearchContexts.occurrence_owner}
\lean{Graph/SearchContexts}{151}{160}{Linear.SearchContexts.selected_fiber}
\lean{Graph/SearchContexts}{163}{173}{Linear.SearchContexts.unique}
\lean{Graph/SearchContexts}{176}{181}{Linear.SearchContexts.present}
\lean{Graph/SearchContexts}{184}{194}{Linear.SearchContexts.only_hole}
\begin{definition}[Focus-free simple search context]\label{term:simple-search-context}
A selected edge carries a specified hole label, the graph is globally
focus-free, and every hole occurrence has that edge as its outer owner.
Because a primitive hole has no child occurrences, this permits exactly one
shallow hole and no others. Opening a slot in a focus-free, hole-free host
supplies these fields. Conversely, closing any such simple context with a
compatible focus-free, hole-free label gives a focus-free, hole-free host.
\end{definition}
The constructor for opening uses the selected-label, focus-freedom, and
occurrence-owner proofs above. The closing theorem pairs relabelling's
focus-freedom theorem with the hole-elimination theorem. Its hypotheses on
the new label are automatically satisfied by finite interaction rules.
\lean{Graph/SearchContexts}{197}{201}{Linear.SearchContexts.Simple}
\lean{Graph/SearchContexts}{204}{205}{Linear.SearchContexts.simple}
\lean{Graph/SearchContexts}{209}{214}{Linear.SearchContexts.close_simple}

\subsection{The explicit-hole correspondence}
\begin{definition}[Explicit-hole pattern search]\label{term:hole-context-search}
Transport a finite fragment's interface to that of the explicitly opened hole.
An explicit-hole match is a graph isomorphism from plugging that transported
fragment into the opened context to the supplied state. Explicit-hole search
chooses one host, slot, and hole name, and one finite interaction rule whose
two fragments match the two states using that same opened context.
\end{definition}
\lean{Graph/HoleContextSearch}{11}{16}{Linear.PatternPlugging.holeFragment}
\lean{Graph/HoleContextSearch}{19}{20}{Linear.PatternPlugging.HoleMatches}
\lean{Graph/HoleContextSearch}{35}{39}{Linear.State.HoleContextSearch}
The discarded-slot-label isomorphism shows that explicit-hole matching and
slot matching are equivalent: compose it with an endpoint match in one
direction and compose its inverse in the other. Hence slot-context search
and explicit-hole search are equivalent. To open a slot witness one may choose
name zero; in the reverse direction any supplied name is retained by the
matching equivalence.
\lean{Graph/HoleContextSearch}{23}{31}{Linear.PatternPlugging.holeMatches_iff}
\lean{Graph/HoleContextSearch}{42}{50}{Linear.State.contextual_iff_holeContext}
\begin{theorem}[Search adequacy for explicit simple contexts]
Erasure-based state search holds exactly when a finite search interaction is
plugged into one explicit focus-free simple hole context, up to graph names.
The same equivalence holds for the independent incidence equations. Recursive
search adds precisely recursive renaming of the resulting state.
\end{theorem}
\begin{proof}
Compose slot-context correspondence with explicit-hole matching equivalence.
Compose once more with the incidence correspondence for the incidence claim.
For recursive search, substitute this equivalence into its existing
literal-step-plus-recursive-renaming characterization.
\end{proof}
\lean{Graph/HoleContextSearch}{54}{56}{Linear.State.search_iff_holeContext}
\lean{Graph/HoleContextSearch}{59}{61}{Linear.State.incidenceSearch_iff_holeContext}
\lean{Graph/HoleContextSearch}{64}{67}{Linear.State.deepSearch_iff_holeContext}

\begin{theorem}[Arbitrary simple contexts are covered]
Take any focus-free simple context, any compatible interaction label, and a
finite corolla rule at that label. If its two fragments, with interface
transport, plug into that original context to match two states, those states
are related by search.
\end{theorem}
\begin{proof}
The rule forces its label to be a variable or operation, so focus freedom and
absence of holes in its label fibre are derived, not added as hypotheses.
Close the original hole with this label. The closing theorem gives a
focus-free, hole-free host. Reverse slot-label transport identifies each
plugging into this host with the supplied plugging into the original context.
Compose with the supplied endpoint isomorphisms to obtain a slot-context search
witness and apply soundness. Thus representing a context by a closed host
and a discarded edge does not exclude an arbitrary simple context.
\end{proof}
\lean{Graph/HoleContextSearch}{72}{95}{Linear.State.arbitrary_simple_context}

\subsection{Regression examples and scope}
The existing terminating value incidence step converts to an explicit-hole
context step with a final target. The variable trigger similarly converts to
a context step ending in rewrite mode. Opening the value edge with name seven
constructs an actual simple context, proves global hole uniqueness, and
exhibits the specified shallow hole.
\lean{Graph/ContextualSearchExamples}{9}{12}{Linear.ContextualSearchExamples.value_contextual}
\lean{Graph/ContextualSearchExamples}{15}{18}{Linear.ContextualSearchExamples.variable_contextual}
\lean{Graph/ContextualSearchExamples}{21}{27}{Linear.ContextualSearchExamples.value_simple_context}
An output-fragment regression checks the actual incidence lists for an
operation with two term arguments and a thunk. The focus source is the fresh
internal vertex and its target is the first output. The operation targets
that fresh vertex, then its unchanged second argument and thunk outputs.
The proof unfolds the finite constructors and evaluates the split tests.
\lean{Graph/ContextualSearchExamples}{31}{41}{Linear.ContextualSearchExamples.output_fragment_incidence}
The cyclic regression defines a unary operation signature and a well-typed
one-edge self-loop with no outer boundary. It reuses the finite self-loop's
incidence laws and proves the new term typing directly. Its input and output
corolla attachments coincide. Both input and output factorization theorems
still apply, providing explicit isomorphisms; no global attachment injectivity
or acyclicity is silently assumed.
\lean{Graph/ContextualSearchExamples}{44}{44}{Linear.ContextualSearchExamples.unaryOps}
\lean{Graph/ContextualSearchExamples}{47}{63}{Linear.ContextualSearchExamples.loop}
\lean{Graph/ContextualSearchExamples}{66}{66}{Linear.ContextualSearchExamples.loop_attachment}
\lean{Graph/ContextualSearchExamples}{69}{71}{Linear.ContextualSearchExamples.loop_input}
\lean{Graph/ContextualSearchExamples}{74}{76}{Linear.ContextualSearchExamples.loop_output}

The former search-adequacy item is removed from Appendix~\ref{app:remaining-work}.
The following section defines the other transition families and assembles the
full UAM. Search adequacy alone does not supply its structural laws or the
concrete specimen and robustness interpretation.
